% =====================================================================
% Appendix E — Legal posture, licensing per record category, FADP
%
% Self-contained legal-posture statement for the OpenCaseLaw resource.
% Section~\ref{sec:corpus} ("Licensing") gives the high-level summary;
% this appendix is the full machine-actionable license matrix plus the
% Swiss-law specifics (copyright posture under URG Art.~5, personal-data
% treatment under FADP). Cited by Sec.~\ref{sec:corpus} and
% Sec.~\ref{sec:limit}.
% =====================================================================

\paragraph{Scope.} This appendix documents the legal posture of the
released resource. It is an operator statement, not a legal opinion;
no Swiss data-protection lawyer has yet reviewed it. A formal
adversarial review is part of the v2.0 release plan.

\paragraph{Record-level license categories.} The corpus is a federation
of materials with heterogeneous upstream provenance. Table~\ref{tab:license_matrix}
gives the per-source breakdown.

\begin{table}[h!]
\centering
\caption{Per-source license category. ``Native'' columns identify the
upstream rights regime; ``Our addition'' columns isolate the
incremental metadata our pipeline produces. Aggregations and derived
graphs combine both layers; see prose below.}
\label{tab:license_matrix}
\begin{tabular}{l l l l}
\toprule
Material & Native rights & Native licence & Our addition \\
\midrule
Federal court decisions (BGer, BVGer, BStGer, \dots) & official act, URG Art.~5 excl. & --- & CC0 \\
Cantonal court decisions & official act, URG Art.~5 excl. & --- & CC0 \\
Federal statutes (Fedlex SR) & official act, URG Art.~5 excl. & --- & CC0 \\
Cantonal statutes (LexWork, SIL, ZH OpenData, \dots) & official act, URG Art.~5 excl. & --- & CC0 \\
Federal Council Botschaften (Fedlex JOLUX) & official act, URG Art.~5 excl. & --- & CC0 \\
entscheidsuche.ch (federation source) & aggregator & CC-BY-4.0 & CC0 (our diff) \\
OnlineKommentar.ch (commentary) & scholarship & CC-BY-4.0 & CC0 (anchor only) \\
OpenLegalCommentary.ch (commentary) & scholarship & CC-BY-SA-4.0 & CC-BY-SA-4.0 (propagated) \\
Citation graph (extracted edges) & derived from URG~5 + our extractor & --- & CC0 \\
Statute graph (extracted edges) & derived from URG~5 + our extractor & --- & CC0 \\
QC ground-truth tables (curator notes) & --- & --- & CC0 \\
Search-quality benchmarks & --- & --- & CC0 \\
Cross-lingual question set & --- & --- & CC0 \\
Code (scrapers, MCP server, builders) & --- & --- & MIT \\
\bottomrule
\end{tabular}
\end{table}

\paragraph{Swiss copyright posture (URG Art.~5).} The Swiss Federal Act
on Copyright (\emph{Urheberrechtsgesetz}, URG / SR 231.1), Art.~5,
excludes from copyright protection ``official acts'' including laws,
decisions, and reports of public authorities. The published text of
court decisions and statutes thus carries no underlying rights restriction.
Our packaging (deterministic identifiers, normalised metadata,
deduplication tables, citation/statute graphs) is an independent work
which we release CC0 to the extent any rights subsist. No portion of
the corpus relies on an asserted database right (\emph{sui generis}
database protection, which Switzerland does not recognise).

\paragraph{Federation and propagation.} Two upstream sources carry
their own licences which we honour:

\begin{itemize}
\item \emph{entscheidsuche.ch} (CC-BY-4.0) is one of our federation
inputs at the per-shard level. Where a decision is sourced from
entscheidsuche we preserve the attribution chain in
\path{decision_source} metadata; downstream redistribution should
honour the upstream attribution requirement.

\item \emph{OpenLegalCommentary.ch} (CC-BY-SA-4.0) commentary text,
when surfaced in the \texttt{get\_commentary} tool, propagates the
share-alike obligation to any derivative work that incorporates the
commentary text verbatim. Our metadata anchors (statute reference,
URL, commentary title, author, license tag) are CC0; the commentary
\emph{body text} retains its CC-BY-SA status. Users who include
verbatim commentary text in a derived corpus must redistribute under
CC-BY-SA-4.0.
\end{itemize}

\paragraph{Personal data treatment (FADP).} Swiss court decisions are
made public by the issuing court; our pipeline does not introduce
personal data that was not already present in the source. We adhere
to the following principles under the revised Federal Act on Data
Protection (FADP, in force since 2023-09-01, RS 235.1):

\begin{itemize}
\item \emph{Source pseudonymisation is inherited unchanged.} Most
federal decisions are pseudonymised by the issuing court (parties
appear as \emph{A.}, \emph{X.\ AG}, etc.); some cantonal decisions
are published in unredacted form by the issuing court itself. We
neither re-pseudonymise material the court chose to publish nor
de-pseudonymise material the court chose to redact.

\item \emph{Takedown procedure.} Removal requests reach us via
\path{team@jonashertner.com} or through the contact link on
\url{https://opencaselaw.ch}. Our standard response is to verify that
the request matches an indexed decision, apply the redaction within
five working days, propagate the change through the nightly publish
pipeline, and append a public note to \path{docs/governance-and-removal-policy.md}
documenting the action (date, decision ID, redaction class).
Verified court-side takedown notifications (where the issuing court
itself withdraws the decision) are honoured the same way without
counter-verification.

\item \emph{No query logging.} The deployed MCP endpoint does not
log query payloads or response payloads. The nginx access log
captures only request user-agent and request volume (Appendix~\ref{app:adoption});
no client matter, document content, or natural-language search
string is retained. Differential-privacy noise is applied to the
public adoption-signal aggregates before publication.

\item \emph{Pro-feature client-side PII redaction.} The Word add-in's
Pro feature set (\emph{Verify}, \emph{Strengthen}, \emph{Find~Support},
\emph{Reflect}) operates on lawyer-drafted text that may contain
client matters. Before any LLM round-trip the add-in applies
client-side structural redaction (e.g.\ AHV numbers, IBANs, AVS/AVS
identifiers, email addresses, phone numbers, dates of birth) and
ships only the redacted payload. The server reapplies the same
redactor as a defence-in-depth layer and rejects requests that still
contain structurally-detectable PII with HTTP~400. We do not log Pro
requests beyond a per-license daily counter (25 calls/day cap).
\end{itemize}

\paragraph{Liability and contact.} The resource is provided
``as-is''. The corpus is a research and analysis substrate, not a
substitute for qualified legal counsel; outputs of any tool built on
the corpus should be verified against the original published
decision. Operational contact for takedown requests, licensing
questions, and incident reports is \path{team@jonashertner.com}; the
governance policy is at
\path{docs/governance-and-removal-policy.md} in the source
repository.

\paragraph{Open items for v2.0.} A formal review by a Swiss
data-protection practitioner is on the v2.0 backlog, as is publication
of a per-record license-tag field in the Parquet export (currently
implicit in the source-court mapping). Aggregating commentary text
into derived models would also require an explicit CC-BY-SA
compliance regime which our current release does not need.
